% GENERATED FILE. Edit canonical lessons, then run npm run generate:books.
% canonical-source-hash: fnv1a64:b6e3bbfed2730f6c
% canonical-lessons: PA-C141-calak, PA-C141-siana, PA-C141-bhola, PA-C141-gol, PA-C141-siddha

\chapter{Cunning, Wise, Innocent, Round, Straight}
\label{ch:pa-cunning-wise-innocent-round-straight}
\hlchaptermodality{141}

\begin{chapteropening}
\textbf{By the end of this chapter:} \emph{I can say cunning, wise, innocent, round and straight.}

Complete the last lesson of chapter 141: I can say cunning, wise, innocent, round and straight.
\end{chapteropening}

\section[\texorpdfstring{\pa{ਚਲਾਕ} (calāk)}{ਚਲਾਕ (calāk)}]{\pa{ਚਲਾਕ} (calāk) — cunning}

\label{lesson:PA-C141-calak}



\begin{quote}
Before the new one: say the Punjabi for light, then the Punjabi for heavy.
\end{quote}

\subsection*{You'll want to know: \pa{ਚਲਾਕ}}
\textbf{\pa{ਚਲਾਕ}} — \emph{calāk} — \textquotedblleft{}cunning\textquotedblright{}.

\begin{grammarlens}[title={What you've built}]
One more word for this chapter.
\end{grammarlens}

\subsection*{Your turn}
\begin{itemize}
\raggedright
  \item \emph{Say it:} \emph{calāk}
  \item \emph{Say it:} \emph{calāk}, once more
  \item \emph{Say it:} say \emph{calāk}
  \item \emph{Recall:} say the Punjabi for light, then the Punjabi for heavy, then say \emph{calāk} again
\end{itemize}

\begin{tcolorbox}[breakable,colback=teal!4,colframe=teal!35!black,title={Before you move on}]
What does \pa{ਚਲਾਕ} mean? (\textquotedblleft{}Cunning\textquotedblright{}.) Say it once more.
\end{tcolorbox}
\section[\texorpdfstring{\pa{ਸਿਆਣਾ} (siāṇā)}{ਸਿਆਣਾ (siāṇā)}]{\pa{ਸਿਆਣਾ} (siāṇā) — wise}

\label{lesson:PA-C141-siana}



\begin{quote}
Before the new one: say the Punjabi for heavy, then the Punjabi for cunning.
\end{quote}

\subsection*{You'll want to know: \pa{ਸਿਆਣਾ}}
\textbf{\pa{ਸਿਆਣਾ}} — \emph{siāṇā} — \textquotedblleft{}wise\textquotedblright{}.

\begin{grammarlens}[title={What you've built}]
One more word for this chapter.
\end{grammarlens}

\subsection*{Your turn}
\begin{itemize}
\raggedright
  \item \emph{Say it:} \emph{siāṇā}
  \item \emph{Say it:} \emph{siāṇā}, once more
  \item \emph{Say it:} say \emph{siāṇā}
  \item \emph{Recall:} say the Punjabi for heavy, then the Punjabi for cunning, then say \emph{siāṇā} again
\end{itemize}

\begin{tcolorbox}[breakable,colback=teal!4,colframe=teal!35!black,title={Before you move on}]
What does \pa{ਸਿਆਣਾ} mean? (\textquotedblleft{}Wise\textquotedblright{}.) Say it once more.
\end{tcolorbox}
\section[\texorpdfstring{\pa{ਭੋਲਾ} (bholā)}{ਭੋਲਾ (bholā)}]{\pa{ਭੋਲਾ} (bholā) — innocent, simple}

\label{lesson:PA-C141-bhola}



\begin{quote}
Before the new one: say the Punjabi for cunning, then the Punjabi for wise.
\end{quote}

\subsection*{You'll want to know: \pa{ਭੋਲਾ}}
\textbf{\pa{ਭੋਲਾ}} — \emph{bholā} — \textquotedblleft{}innocent, simple\textquotedblright{}.

\begin{grammarlens}[title={What you've built}]
One more word for this chapter.
\end{grammarlens}

\subsection*{Your turn}
\begin{itemize}
\raggedright
  \item \emph{Say it:} \emph{bholā}
  \item \emph{Say it:} \emph{bholā}, once more
  \item \emph{Say it:} say \emph{bholā}
  \item \emph{Recall:} say the Punjabi for cunning, then the Punjabi for wise, then say \emph{bholā} again
\end{itemize}

\begin{tcolorbox}[breakable,colback=teal!4,colframe=teal!35!black,title={Before you move on}]
What does \pa{ਭੋਲਾ} mean? (\textquotedblleft{}Innocent, simple\textquotedblright{}.) Say it once more.
\end{tcolorbox}
\section[\texorpdfstring{\pa{ਗੋਲ} (gol)}{ਗੋਲ (gol)}]{\pa{ਗੋਲ} (gol) — round}

\label{lesson:PA-C141-gol}



\begin{quote}
Before the new one: say the Punjabi for wise, then the Punjabi for innocent.
\end{quote}

\subsection*{You'll want to know: \pa{ਗੋਲ}}
\textbf{\pa{ਗੋਲ}} — \emph{gol} — \textquotedblleft{}round\textquotedblright{}.

\begin{grammarlens}[title={What you've built}]
One more word for this chapter.
\end{grammarlens}

\subsection*{Your turn}
\begin{itemize}
\raggedright
  \item \emph{Say it:} \emph{gol}
  \item \emph{Say it:} \emph{gol}, once more
  \item \emph{Say it:} say \emph{gol}
  \item \emph{Recall:} say the Punjabi for wise, then the Punjabi for innocent, then say \emph{gol} again
\end{itemize}

\begin{tcolorbox}[breakable,colback=teal!4,colframe=teal!35!black,title={Before you move on}]
What does \pa{ਗੋਲ} mean? (\textquotedblleft{}Round\textquotedblright{}.) Say it once more.
\end{tcolorbox}
\section[\texorpdfstring{\pa{ਸਿੱਧਾ} (siddhā)}{ਸਿੱਧਾ (siddhā)}]{\pa{ਸਿੱਧਾ} (siddhā) — straight}

\label{lesson:PA-C141-siddha}



\begin{quote}
Before the new one: say the Punjabi for innocent, then the Punjabi for round.
\end{quote}

\subsection*{You'll want to know: \pa{ਸਿੱਧਾ}}
\textbf{\pa{ਸਿੱਧਾ}} — \emph{siddhā} — \textquotedblleft{}straight\textquotedblright{}.

\begin{grammarlens}[title={What you've built}]
One more word for this chapter.
\end{grammarlens}

\subsection*{Your turn}
\begin{itemize}
\raggedright
  \item \emph{Say it:} \emph{siddhā}
  \item \emph{Say it:} \emph{siddhā}, once more
  \item \emph{Say it:} say \emph{siddhā}
  \item \emph{Recall:} say the Punjabi for innocent, then the Punjabi for round, then say \emph{siddhā} again
\end{itemize}

\begin{tcolorbox}[breakable,colback=teal!4,colframe=teal!35!black,title={Before you move on}]
What does \pa{ਸਿੱਧਾ} mean? (\textquotedblleft{}Straight\textquotedblright{}.) Say it once more.
\end{tcolorbox}
